% =====================================================================
\chapter{Soluciones de los ejercicios}
\label{ap:soluciones}
% =====================================================================

Cada solución se presenta en cinco partes: datos, planteamiento, desarrollo, resultado y comprobación. En los pasos intermedios se redondea a cuatro cifras significativas salvo que se indique otra cosa; los resultados exactos se dan como fracción cuando existe. Los valores de $\Phi$ proceden del \cref{tab:normal} o de un cálculo numérico.

% ---------------------------------------------------------------------
\section{Tema 1}

\begin{solucion}{eje:1-inversa}
\datos se pide un ejemplo con valores plausibles.
\planteamiento por el \cref{teo:bayes}, $\PP(A\mid B)=\PP(B\mid A)\,\PP(A)/\PP(B)$: las dos condicionadas difieren en el factor $\PP(A)/\PP(B)$.
\pasos
\begin{enumerate}[label=\arabic*.]
  \item Sean $A=\{\text{llueve}\}$ y $B=\{\text{la calle está mojada}\}$.
  \item Si llueve, la calle se moja: $\PP(B\mid A)\approx 1$.
  \item La calle también se moja por el riego o la limpieza municipal; se toma $\PP(A)=0{,}1$ y $\PP(B)=0{,}4$.
  \item $\PP(A\mid B)=1\cdot 0{,}1/0{,}4=0{,}25$.
\end{enumerate}
\resultado \textbf{$\PP(B\mid A)\approx 1$, pero $\PP(A\mid B)\approx 0{,}25$}; las relaciona el cociente $\PP(A)/\PP(B)=0{,}25$.
\comprobacion ambas factorizaciones dan la misma probabilidad conjunta: $\PP(B\mid A)\PP(A)=1\cdot 0{,}1=0{,}1$ y $\PP(A\mid B)\PP(B)=0{,}25\cdot 0{,}4=0{,}1$.
\end{solucion}

\begin{solucion}{eje:1-disjuntos}
\datos $A\cap B=\varnothing$, $\PP(A)>0$, $\PP(B)>0$.
\planteamiento comparar $\PP(A\cap B)$ con $\PP(A)\PP(B)$ (\cref{def:indep}).
\pasos
\begin{enumerate}[label=\arabic*.]
  \item Como $A\cap B=\varnothing$, $\PP(A\cap B)=\PP(\varnothing)=0$ por la \cref{prop:axiomas}(i).
  \item $\PP(A)\PP(B)>0$ por ser producto de dos números positivos.
  \item Luego $\PP(A\cap B)\ne\PP(A)\PP(B)$ y los sucesos no son independientes.
\end{enumerate}
\resultado \textbf{Dos sucesos incompatibles con probabilidad positiva son dependientes.} Intuitivamente, saber que ha ocurrido $A$ da la máxima información sobre $B$: que no ha ocurrido, pues $\PP(B\mid A)=0\ne\PP(B)$. $\blacksquare$
\comprobacion con un dado, $A=\{1\}$ y $B=\{2\}$: $\PP(A\cap B)=0$ y $\PP(A)\PP(B)=1/36$.
\end{solucion}

\begin{solucion}{eje:1-detector}
\datos $\PP(G)=0{,}01$ (generado); $\PP(M\mid G)=0{,}99$ (marcado); $\PP(M\mid G^{c})=0{,}05$.
\planteamiento regla de Bayes con la probabilidad total en el denominador (\cref{teo:total,teo:bayes}).
\pasos
\begin{enumerate}[label=\arabic*.]
  \item Verdaderos positivos: $\PP(M\mid G)\PP(G)=0{,}99\cdot 0{,}01=0{,}0099$.
  \item Falsos positivos: $\PP(M\mid G^{c})\PP(G^{c})=0{,}05\cdot 0{,}99=0{,}0495$.
  \item Evidencia: $\PP(M)=0{,}0099+0{,}0495=0{,}0594$.
  \item $\PP(G\mid M)=0{,}0099/0{,}0594=1/6\approx 0{,}1667$.
\end{enumerate}
\resultado \textbf{$\PP(G\mid M)=1/6\approx 0{,}167$}: cinco de cada seis alarmas son falsas.
\comprobacion con $10\,000$ textos hay 100 generados (99 marcados) y $9900$ humanos (495 marcados); $99/(99+495)=99/594=1/6$.
\end{solucion}

\begin{solucion}{eje:1-cribado}
\datos $\PP(E)=0{,}005$; $\PP(+\mid E^{c})=0{,}05$; $\PP(-\mid E)=0{,}02$, luego $\PP(+\mid E)=0{,}98$.
\planteamiento \cref{prop:odds}: razón a posteriori $=$ factor de Bayes $\times$ razón a priori; después $\PP=O/(1+O)$.
\pasos
\begin{enumerate}[label=\arabic*.]
  \item Razón a priori: $O(E)=0{,}005/0{,}995=1/199\approx 0{,}005025$.
  \item Factor de Bayes: $0{,}98/0{,}05=19{,}6$.
  \item Razón a posteriori: $O(E\mid +)=19{,}6/199=98/995\approx 0{,}09849$.
  \item Probabilidad a posteriori: $0{,}09849/1{,}09849\approx 0{,}08966$.
\end{enumerate}
\resultado \textbf{razón a priori $1/199$; factor de Bayes $19{,}6$; razón a posteriori $\approx 0{,}0985$; $\PP(E\mid +)\approx 0{,}0897$.}
\comprobacion por la vía directa, $\dfrac{0{,}98\cdot 0{,}005}{0{,}98\cdot 0{,}005+0{,}05\cdot 0{,}995}=\dfrac{0{,}0049}{0{,}05465}=0{,}08966$. La prueba es informativa (multiplica la razón por casi~20), pero la razón a priori es tan baja que un positivo sigue siendo probablemente falso.
\end{solucion}

\begin{solucion}{eje:1-monedas}
\datos $\PP(J)=\PP(J^{c})=1/2$ ($J$: moneda equilibrada); $\PP(\text{cara}\mid J)=1/2$; $\PP(\text{cara}\mid J^{c})=1/4$; dada la moneda, los lanzamientos son independientes.
\planteamiento (a) regla de Bayes; (b) probabilidad total para $\PP(C_1)$, $\PP(C_2)$ y $\PP(C_1\cap C_2)$, donde $C_i$ es «cara en el lanzamiento $i$».
\pasos
\begin{enumerate}[label=\arabic*.]
  \item $\PP(C_1\cap C_2\mid J)=(1/2)^{2}=1/4$ y $\PP(C_1\cap C_2\mid J^{c})=(1/4)^{2}=1/16$.
  \item $\PP(C_1\cap C_2)=\frac12\cdot\frac14+\frac12\cdot\frac1{16}=\frac18+\frac1{32}=\frac{5}{32}$.
  \item (a) $\PP(J\mid C_1\cap C_2)=\dfrac{1/8}{5/32}=\dfrac{4}{5}$.
  \item (b) $\PP(C_1)=\frac12\cdot\frac12+\frac12\cdot\frac14=\frac38$, y por simetría $\PP(C_2)=\frac38$.
  \item $\PP(C_1)\PP(C_2)=\frac{9}{64}=0{,}1406\ne\frac{5}{32}=0{,}1563=\PP(C_1\cap C_2)$.
\end{enumerate}
\resultado \textbf{(a) $4/5$; (b) no son independientes}, aunque sí lo son condicionalmente a la moneda: una cara en el primer lanzamiento aumenta la probabilidad de la moneda equilibrada y, con ella, la de cara en el segundo, $\PP(C_2\mid C_1)=(5/32)/(3/8)=5/12>3/8$.
\comprobacion la probabilidad a priori de la moneda equilibrada era $1/2$; dos caras la elevan a $4/5$, coherente con que esa moneda produce caras con mayor frecuencia.
\end{solucion}

\begin{solucion}{eje:1-complementos}
\datos $\PP(A\cap B)=\PP(A)\PP(B)$.
\planteamiento descomponer $A$ en una unión disjunta y usar la aditividad; después aplicar el resultado dos veces.
\pasos
\begin{enumerate}[label=\arabic*.]
  \item $A=(A\cap B)\cup(A\cap B^{c})$ con unión disjunta, luego $\PP(A\cap B^{c})=\PP(A)-\PP(A\cap B)$.
  \item Por la hipótesis, $\PP(A\cap B^{c})=\PP(A)-\PP(A)\PP(B)=\PP(A)\bigl(1-\PP(B)\bigr)=\PP(A)\PP(B^{c})$, por la \cref{prop:axiomas}(i). Así, $A$ y $B^{c}$ son independientes.
  \item La relación de independencia es simétrica. Aplicando el paso~2 al par independiente $(B^{c},A)$, se obtiene que $B^{c}$ y $A^{c}$ son independientes.
\end{enumerate}
\resultado \textbf{Si $A$ y $B$ son independientes, también lo son $A$ y $B^{c}$, y $A^{c}$ y $B^{c}$.} $\blacksquare$
\comprobacion con dos monedas equilibradas, $A=\{\text{cara en la 1.ª}\}$ y $B=\{\text{cara en la 2.ª}\}$: $\PP(A^{c}\cap B^{c})=1/4=\PP(A^{c})\PP(B^{c})$.
\end{solucion}

% ---------------------------------------------------------------------
\section{Tema 2}

\begin{solucion}{eje:2-categorica}
\datos la salida de un LLM es un índice de token en $\{1,\dots,K\}$.
\planteamiento comparar lo que supone cada distribución sobre los valores que modela.
\pasos
\begin{enumerate}[label=\arabic*.]
  \item Los tokens son categorías sin orden: los índices son etiquetas arbitrarias, y el token~517 no se parece más al~518 que al~30\,000.
  \item Una normal vive en $\R$: asignaría densidad a valores no enteros y supondría que índices próximos son parecidos, porque su densidad es una función suave del índice.
  \item Una normal es unimodal: no puede repartir la masa entre dos tokens de índices lejanos («gato» y «perro») sin ponerla también en los intermedios.
  \item La categórica asigna una probabilidad independiente a cada token, con $p_k\ge 0$ y $\sum_k p_k=1$, que es justo lo que produce la softmax.
\end{enumerate}
\resultado \textbf{la categórica, porque los tokens son categorías discretas sin orden y la distribución del token siguiente puede ser multimodal.}
\comprobacion con $K=50\,000$, la categórica tiene $49\,999$ parámetros libres y la normal, 2: ninguna familia de dos parámetros puede representar preferencias arbitrarias entre $50\,000$ opciones.
\end{solucion}

\begin{solucion}{eje:2-dados}
\datos $S=X_1+X_2$, dados equilibrados independientes; por el \cref{ej:tresvalores}, $\E[X_i]=7/2$ y $\Var(X_i)=35/12$.
\planteamiento linealidad (\cref{prop:linealidad}), suma de varianzas de independientes (\cref{prop:suma}) y recuento de casos favorables.
\pasos
\begin{enumerate}[label=\arabic*.]
  \item $\E[S]=7/2+7/2=7$.
  \item $\Var(S)=35/12+35/12=35/6\approx 5{,}833$.
  \item Pares con suma~7: $(1,6),(2,5),(3,4),(4,3),(5,2),(6,1)$, seis de los 36 equiprobables.
  \item $\PP(S=7)=6/36=1/6\approx 0{,}1667$.
\end{enumerate}
\resultado \textbf{$\E[S]=7$; $\Var(S)=35/6\approx 5{,}833$; $\PP(S=7)=1/6$.}
\comprobacion $\E[X_i^{2}]-\E[X_i]^{2}=91/6-49/4=182/12-147/12=35/12$. Además, $7$ es la suma más probable y coincide con $\E[S]$, por la simetría de la distribución de $S$.
\end{solucion}

\begin{solucion}{eje:2-softmax}
\datos $z=(2,1,0,-1)$; $u=0{,}5$ y $u=0{,}9$; tokens numerados de~0 a~3.
\planteamiento $p_k=e^{z_k}/\sum_j e^{z_j}$; se emite el primer token cuya acumulada sea $\ge u$ (\cref{prop:inversa-discreta}).
\pasos
\begin{enumerate}[label=\arabic*.]
  \item $e^{z}=(7{,}389;\ 2{,}718;\ 1;\ 0{,}3679)$.
  \item Suma: $7{,}389+2{,}718+1+0{,}3679=11{,}475$ (se conservan cinco cifras para que los cocientes tengan tres decimales correctos).
  \item $p=(7{,}389;\ 2{,}718;\ 1;\ 0{,}3679)/11{,}475=(0{,}644;\ 0{,}237;\ 0{,}087;\ 0{,}032)$.
  \item Acumuladas: $F=(0{,}644;\ 0{,}881;\ 0{,}968;\ 1)$.
  \item $u=0{,}5\le 0{,}644$: se emite el token~0.
  \item $0{,}881<0{,}9\le 0{,}968$: con $u=0{,}9$ se emite el token~2.
\end{enumerate}
\resultado \textbf{$p=(0{,}644;\,0{,}237;\,0{,}087;\,0{,}032)$; $F=(0{,}644;\,0{,}881;\,0{,}968;\,1)$; token~0 con $u=0{,}5$ y token~2 con $u=0{,}9$.}
\comprobacion $0{,}644+0{,}237+0{,}087+0{,}032=1{,}000$.
\end{solucion}

\begin{solucion}{eje:2-dados4}
\datos $\PP(i,j)=c\,(i+j)$ para $i,j\in\{1,2,3,4\}$.
\planteamiento normalización para hallar $c$; después, sumar $c(i+j)$ sobre los pares de cada suceso; para (d), \cref{def:esperanza}.
\pasos
\begin{enumerate}[label=\arabic*.]
  \item $\sum_{i,j}(i+j)=\sum_{i,j}i+\sum_{i,j}j=4\cdot 10+4\cdot 10=80$, luego $c=1/80$.
  \item (b) Sumas pares: 2 (1 par), 4 (3 pares), 6 (3 pares), 8 (1 par). Peso: $2\cdot 1+4\cdot 3+6\cdot 3+8\cdot 1=40$; $\PP=40/80=1/2$.
  \item (c) Pares $(2,3)$ y $(3,2)$, de peso $5$ cada uno: $\PP=10/80=1/8$.
  \item (d) $\E[S]=\sum_{i,j}(i+j)\cdot\frac{i+j}{80}=\frac1{80}\sum_s s^{2}\,n_s$, con $n_s$ el número de pares de suma $s$: $4\cdot 1+9\cdot 2+16\cdot 3+25\cdot 4+36\cdot 3+49\cdot 2+64\cdot 1=440$.
  \item $\E[S]=440/80=11/2=5{,}5$.
\end{enumerate}
\resultado \textbf{$c=1/80$; $\PP(\text{par})=1/2$; $\PP(\{2,3\})=1/8$; $\E[S]=5{,}5$.}
\comprobacion con dados equilibrados de cuatro caras $\E[S]=5$; aquí los pesos favorecen las sumas altas, de modo que $5{,}5>5$ es coherente.
\end{solucion}

\begin{solucion}{eje:2-binomial}
\datos $X\sim\Bin(10;\,0{,}8)$.
\planteamiento $\E[X]=np$, $\Var(X)=np(1-p)$ y la función de masa binomial (\cref{def:discretas}).
\pasos
\begin{enumerate}[label=\arabic*.]
  \item $\E[X]=10\cdot 0{,}8=8$.
  \item $\Var(X)=10\cdot 0{,}8\cdot 0{,}2=1{,}6$.
  \item $\PP(X=10)=0{,}8^{10}=0{,}1074$.
  \item $\PP(X=9)=\binom{10}{9}0{,}8^{9}\cdot 0{,}2=10\cdot 0{,}1342\cdot 0{,}2=0{,}2684$.
  \item $\PP(X\ge 9)=0{,}2684+0{,}1074=0{,}3758$.
\end{enumerate}
\resultado \textbf{$\E[X]=8$; $\Var(X)=1{,}6$; $\PP(X=10)\approx 0{,}107$; $\PP(X\ge 9)\approx 0{,}376$.}
\comprobacion $\PP(X=9)/\PP(X=10)=10\cdot 0{,}2/0{,}8=2{,}5$ y $0{,}2684/0{,}1074=2{,}50$.
\end{solucion}

\begin{solucion}{eje:2-covarianza}
\datos (a) $X$, $Y$ con varianza finita. (b) $\PP(X=-1)=\PP(X=0)=\PP(X=1)=1/3$, $Y=X^{2}$.
\planteamiento (a) desarrollar la definición de varianza; (b) calcular la covarianza y buscar un par $(x,y)$ en que la conjunta no factorice.
\pasos
\begin{enumerate}[label=\arabic*.]
  \item (a) Sean $\mu_X=\E[X]$ y $\mu_Y=\E[Y]$. Por linealidad, $\E[X+Y]=\mu_X+\mu_Y$.
  \item $(X+Y-\mu_X-\mu_Y)^{2}=(X-\mu_X)^{2}+(Y-\mu_Y)^{2}+2(X-\mu_X)(Y-\mu_Y)$.
  \item Tomando esperanzas y usando de nuevo la linealidad: $\Var(X+Y)=\Var(X)+\Var(Y)+2\Cov(X,Y)$.
  \item $\Cov(X,Y)=\E[XY]-\mu_X\E[Y]-\mu_Y\E[X]+\mu_X\mu_Y=\E[XY]-\mu_X\mu_Y$, de nuevo por linealidad.
  \item Si son independientes, $\E[XY]=\sum_{x,y}xy\,p_X(x)p_Y(y)=\mu_X\mu_Y$ y la covarianza se anula.
  \item (b) $\E[X]=(-1+0+1)/3=0$ y $\E[XY]=\E[X^{3}]=(-1+0+1)/3=0$, luego $\Cov(X,Y)=0$.
  \item $\PP(X=0,Y=0)=\PP(X=0)=1/3$, mientras que $\PP(X=0)\PP(Y=0)=\frac13\cdot\frac13=\frac19$.
\end{enumerate}
\resultado \textbf{(a) la fórmula de la \cref{prop:suma}; (b) $X$ e $Y$ son incorreladas pero dependientes}: $Y$ es función de $X$. $\blacksquare$
\comprobacion la covarianza solo detecta dependencia lineal; aquí la dependencia es cuadrática y simétrica, y las contribuciones de $x=-1$ y $x=1$ se cancelan.
\end{solucion}

% ---------------------------------------------------------------------
\section{Tema 3}

\begin{solucion}{eje:3-tabla}
\datos $X\sim\N(10,4)$, luego $\mu=10$ y $\sigma=2$.
\planteamiento estandarizar con la \cref{prop:afin} y usar el \cref{tab:normal} y la simetría.
\pasos
\begin{enumerate}[label=\arabic*.]
  \item $8\le X\le 12$ equivale a $\abs{X-10}\le 2=1\cdot\sigma$: probabilidad $0{,}683$.
  \item $X>14$ equivale a $Z=(X-10)/2>2$.
  \item $\PP(\abs Z\le 2)=0{,}954$, luego $\PP(\abs Z>2)=0{,}046$ y, por simetría, $\PP(Z>2)=0{,}046/2=0{,}023$.
\end{enumerate}
\resultado \textbf{$\PP(8\le X\le 12)\approx 0{,}683$; $\PP(X>14)\approx 0{,}023$.}
\comprobacion numéricamente, $1-\Phi(2)=0{,}02275$.
\end{solucion}

\begin{solucion}{eje:3-espera}
\datos $X\sim\Expo(0{,}5)$, en minutos.
\planteamiento $\E[X]=1/\lambda$, $\Var(X)=1/\lambda^{2}$, $\PP(X>x)=e^{-\lambda x}$ y la \cref{prop:memoria}.
\pasos
\begin{enumerate}[label=\arabic*.]
  \item (a) $\E[X]=1/0{,}5=2$ minutos; $\sigma_X=\sqrt{1/0{,}25}=2$ minutos.
  \item (b) $\PP(X>2)=e^{-0{,}5\cdot 2}=e^{-1}=0{,}3679$.
  \item (c) $\PP(X>5\mid X>3)=\PP(X>3+2\mid X>3)=\PP(X>2)=e^{-1}$.
\end{enumerate}
\resultado \textbf{(a) media y desviación típica de 2 minutos; (b) y (c) $e^{-1}\approx 0{,}368$.}
\comprobacion directamente, $\PP(X>5)/\PP(X>3)=e^{-2{,}5}/e^{-1{,}5}=e^{-1}$: los 3 minutos ya esperados no cambian la distribución del tiempo restante.
\end{solucion}

\begin{solucion}{eje:3-cuadrado}
\datos $U\sim\Unif(-1,1)$, de densidad $1/2$ en $[-1,1]$.
\planteamiento \cref{prop:uniforme} para la media y la varianza; integral directa para $\E[U^{4}]$; función de distribución de $U^{2}$ a partir de la de $U$.
\pasos
\begin{enumerate}[label=\arabic*.]
  \item $\E[U]=(-1+1)/2=0$ y $\Var(U)=2^{2}/12=1/3$.
  \item $\E[U^{4}]=\frac12\int_{-1}^{1}u^{4}\,du=\frac12\cdot\frac{2}{5}=\frac15$.
  \item Para $t\in[0,1]$: $\PP(U^{2}\le t)=\PP(-\sqrt t\le U\le\sqrt t)=\frac{2\sqrt t}{2}=\sqrt t$. Vale $0$ si $t<0$ y $1$ si $t>1$.
  \item Derivando, la densidad de $U^{2}$ es $\frac{1}{2\sqrt t}$ en $(0,1]$.
\end{enumerate}
\resultado \textbf{$\E[U]=0$, $\Var(U)=1/3$, $\E[U^{4}]=1/5$; $F_{U^{2}}(t)=\sqrt t$ y $f_{U^{2}}(t)=1/(2\sqrt t)$ en $(0,1]$: $U^{2}$ no es uniforme}, pues su densidad no es constante y acumula masa cerca de~0.
\comprobacion $\E[U^{2}]=\int_0^1 t\cdot\frac{1}{2\sqrt t}\,dt=\frac12\cdot\frac23=\frac13=\Var(U)$, como debe ser porque $\E[U]=0$.
\end{solucion}

\begin{solucion}{eje:3-init}
\datos $n=4096$, $\Var(x)=1$; en (c), $\sigma_w=0{,}1$ y $n=256$.
\planteamiento \cref{prop:neurona}: $\Var(y)=n\Var(w)\Var(x)$; regla de fan-in $\Var(w)=1/n$; regla de He $\Var(w)=2/n$.
\pasos
\begin{enumerate}[label=\arabic*.]
  \item (a) $\sigma_w=1/\sqrt{4096}=1/64=0{,}01563$.
  \item (b) $\sigma_w=\sqrt{2/4096}=\sqrt2/64=1{,}414/64=0{,}02210$.
  \item (c) $\Var(y)=256\cdot 0{,}1^{2}\cdot 1=2{,}56$, de modo que la desviación típica se multiplica por $\sqrt{2{,}56}=1{,}6$ en cada capa.
  \item Tras diez capas, el factor es $1{,}6^{10}\approx 110$.
\end{enumerate}
\resultado \textbf{(a) $\sigma_w=1/64\approx 0{,}0156$; (b) $\sigma_w\approx 0{,}0221$; (c) $\Var(y)=2{,}56$: la señal crece un factor $1{,}6$ por capa, unas 110 veces en diez capas.}
\comprobacion la desviación correcta para $n=256$ es $1/\sqrt{256}=0{,}0625<0{,}1$; con ella, $256\cdot 0{,}0625^{2}=1$.
\end{solucion}

\begin{solucion}{eje:3-exponencial}
\datos $U\sim\Unif(0,1)$; $F(x)=1-e^{-\lambda x}$ para $x\ge 0$.
\planteamiento aplicar el \cref{teo:inversa}: invertir $F$ y comprobar la función de distribución resultante.
\pasos
\begin{enumerate}[label=\arabic*.]
  \item $F$ es continua y estrictamente creciente en $[0,\infty)$, con valores en $[0,1)$.
  \item Despejando $u=1-e^{-\lambda x}$: $F^{-1}(u)=-\ln(1-u)/\lambda$.
  \item Sea $X=-\ln(1-U)/\lambda$. Para $x\ge 0$: $X\le x\iff\ln(1-U)\ge-\lambda x\iff 1-U\ge e^{-\lambda x}\iff U\le 1-e^{-\lambda x}$.
  \item Luego $\PP(X\le x)=1-e^{-\lambda x}=F(x)$, y $X\sim\Expo(\lambda)$.
  \item Si $U\sim\Unif(0,1)$, también $1-U\sim\Unif(0,1)$, pues $\PP(1-U\le c)=\PP(U\ge 1-c)=c$. Sustituyendo $1-U$ por $U$, $-\ln(U)/\lambda\sim\Expo(\lambda)$.
\end{enumerate}
\resultado \textbf{$X=-\ln(1-U)/\lambda$, o equivalentemente $-\ln(U)/\lambda$, sigue una $\Expo(\lambda)$.} $\blacksquare$
\comprobacion con $\lambda=1$ y $u=1-e^{-1}=0{,}6321$, se obtiene $x=1$, la media de la $\Expo(1)$; la mitad de las muestras cae por debajo de $\ln 2/\lambda$, la mediana.
\end{solucion}

\begin{solucion}{eje:3-he}
\datos $x$ con densidad $f$ simétrica, $f(-t)=f(t)$, y $\Var(x)=v$; $r=\max(0,x)$; $w_j$ con media~0, independientes entre sí y de los $r_j$, que son independientes entre sí.
\planteamiento (a) integrar usando la simetría; (b) varianza de un producto de independientes y \cref{prop:suma}; (c) sustituir.
\pasos
\begin{enumerate}[label=\arabic*.]
  \item (a) Por simetría, $\E[x]=0$, luego $\E[x^{2}]=v$.
  \item $\E[r^{2}]=\int_0^{\infty}t^{2}f(t)\,dt$, porque $r=0$ cuando $x<0$.
  \item La función $t^{2}f(t)$ es par, así que $\int_0^{\infty}t^{2}f(t)\,dt=\frac12\int_{-\infty}^{\infty}t^{2}f(t)\,dt=\frac v2$.
  \item (b) $\E[w_jr_j]=\E[w_j]\E[r_j]=0$ y $\Var(w_jr_j)=\E[w_j^{2}]\E[r_j^{2}]-0=\Var(w)\,\E[r^{2}]$. No hace falta $\E[r_j]=0$: basta con que $\E[w_j]=0$.
  \item Los productos $w_jr_j$ son independientes, luego $\Var(y)=\sum_j\Var(w_jr_j)=n\Var(w)\,\E[r^{2}]$.
  \item (c) Con $\Var(w)=2/n$: $\Var(y)=n\cdot\frac2n\cdot\frac v2=v$.
\end{enumerate}
\resultado \textbf{$\E[r^{2}]=v/2$, $\Var(y)=n\Var(w)\E[r^{2}]$, y la regla de He $\Var(w)=2/n$ conserva la varianza $v$.} $\blacksquare$
\comprobacion con $n=512$, $\sigma_w=\sqrt{2/512}=0{,}0625$, el valor dado en el texto.
\end{solucion}

% ---------------------------------------------------------------------
\section{Tema 4}

\begin{solucion}{eje:4-ceros}
\datos conteos $N=(3,1,0)$, $K=3$, $N_{\mathrm{tot}}=4$.
\planteamiento EMV $N_j/N_{\mathrm{tot}}$ frente a media a posteriori con $\alpha=1$ (\cref{cor:laplace}).
\pasos
\begin{enumerate}[label=\arabic*.]
  \item El EMV asigna $0/4=0$ al tercer valor. Si ese valor aparece después, la verosimilitud es nula y la entropía cruzada, infinita.
  \item Con pocos datos, muchos valores posibles no se observan por azar; el EMV confunde «no observado» con «imposible».
  \item Una a priori $\Dir(\alpha,\dots,\alpha)$ añade $\alpha$ pseudoconteos a cada valor, de modo que ninguna probabilidad estimada es nula.
  \item Con $\alpha=1$: $(0+1)/(4+3)=1/7\approx 0{,}1429$.
\end{enumerate}
\resultado \textbf{el suavizado de Laplace asigna $1/7$ al tercer valor}, frente al~$0$ del EMV.
\comprobacion las tres estimaciones, $(4/7,\,2/7,\,1/7)$, suman~1 (\cref{tab:estimadores}).
\end{solucion}

\begin{solucion}{eje:4-datos}
\datos $x=(1,3,5,7)$, $n=4$.
\planteamiento \cref{teo:emv-normal} y \cref{prop:sesgo}.
\pasos
\begin{enumerate}[label=\arabic*.]
  \item $\hat\mu=(1+3+5+7)/4=16/4=4$.
  \item Desviaciones: $-3,-1,1,3$; cuadrados: $9,1,1,9$; suma $20$.
  \item $\hat\sigma^{2}=20/4=5$.
  \item $s^{2}=20/3\approx 6{,}667$.
\end{enumerate}
\resultado \textbf{$\hat\mu=4$; $\hat\sigma^{2}=5$; $s^{2}=20/3\approx 6{,}667$.}
\comprobacion por la \cref{prop:varianza}(i) aplicada a la distribución empírica, $(1+9+25+49)/4-4^{2}=21-16=5$.
\end{solucion}

\begin{solucion}{eje:4-moneda}
\datos $k=7$, $n=10$.
\planteamiento \cref{ej:emv-bernoulli}; Laplace es la media a posteriori con a priori $\BetaD(1,1)$ (\cref{teo:beta}).
\pasos
\begin{enumerate}[label=\arabic*.]
  \item $\hat p=7/10=0{,}7$.
  \item $\ln 0{,}7=-0{,}3567$ y $\ln 0{,}3=-1{,}204$.
  \item $\ell(0{,}7)=7\cdot(-0{,}3567)+3\cdot(-1{,}204)=-2{,}497-3{,}612=-6{,}109$.
  \item $\ell(0{,}5)=10\cdot(-0{,}6931)=-6{,}931$.
  \item Laplace: $(7+1)/(10+2)=8/12=2/3\approx 0{,}6667$.
\end{enumerate}
\resultado \textbf{$\hat p=0{,}7$; $\ell(0{,}7)\approx-6{,}109>\ell(0{,}5)\approx-6{,}931$; Laplace, $2/3\approx 0{,}667$.} La estimación bayesiana queda entre la media a priori, $1/2$, y el EMV, $0{,}7$: con pocos datos, la a priori atrae la estimación hacia el centro.
\comprobacion $\ell'(0{,}7)=7/0{,}7-3/0{,}3=10-10=0$.
\end{solucion}

\begin{solucion}{eje:4-u2}
\datos 49 caras y 31 cruces en 80 lanzamientos; duraciones $2,3,1,3,4$.
\planteamiento comparar $\ell(p)=49\ln p+31\ln(1-p)$ (el coeficiente binomial es común a todas las hipótesis); para la exponencial, $f(x)=\lambda e^{-\lambda x}$.
\pasos
\begin{enumerate}[label=\arabic*.]
  \item (a) $p=1/3$: $\ell=49\cdot(-1{,}099)+31\cdot(-0{,}4055)=-53{,}83-12{,}57=-66{,}40$.
  \item $p=1/2$: $\ell=80\cdot(-0{,}6931)=-55{,}45$.
  \item $p=2/3$: $\ell=49\cdot(-0{,}4055)+31\cdot(-1{,}099)=-19{,}87-34{,}06=-53{,}92$, el mayor de los tres.
  \item (b) $\ell'(p)=49/p-31/(1-p)=0\Rightarrow 49(1-p)=31p\Rightarrow\hat p=49/80=0{,}6125$.
  \item (c) $\ell(\lambda)=\sum_i(\ln\lambda-\lambda x_i)=5\ln\lambda-13\lambda$, pues $\sum_i x_i=13$.
  \item $\ell'(\lambda)=5/\lambda-13=0\Rightarrow\hat\lambda=5/13\approx 0{,}3846$; $\ell''(\lambda)=-5/\lambda^{2}<0$, máximo.
\end{enumerate}
\resultado \textbf{(a) la moneda de $p=2/3$; (b) $\hat p=49/80=0{,}6125$; (c) $\hat\lambda=5/13\approx 0{,}385$ por año} (vida media estimada, $2{,}6$ años).
\comprobacion $\ell(5/13)=5\ln 0{,}3846-5=-9{,}778$, mayor que $\ell(0{,}3)=-9{,}920$ y que $\ell(0{,}5)=-9{,}966$.
\end{solucion}

\begin{solucion}{eje:4-u3}
\datos a priori $1/3$ para cada moneda; dato: una cara en 5 lanzamientos.
\planteamiento verosimilitud binomial $5p(1-p)^{4}$; tabla de Bayes; probabilidad predictiva por probabilidad total, $\sum_i\PP(\text{cara}\mid i)\PP(i\mid\text{datos})$.
\pasos
\begin{enumerate}[label=\arabic*.]
  \item Verosimilitudes: $p=0{,}5$: $5\cdot 0{,}5\cdot 0{,}5^{4}=0{,}1563$; $p=0{,}6$: $5\cdot 0{,}6\cdot 0{,}4^{4}=0{,}0768$; $p=0{,}9$: $5\cdot 0{,}9\cdot 0{,}1^{4}=0{,}00045$.
  \item Productos por la a priori ($\div 3$): $0{,}05208$; $0{,}0256$; $0{,}00015$.
  \item Evidencia: $0{,}05208+0{,}0256+0{,}00015=0{,}07783$.
  \item A posteriori: $0{,}6692$; $0{,}3289$; $0{,}001927$.
  \item Predictiva a priori: $(0{,}5+0{,}6+0{,}9)/3=2/3\approx 0{,}6667$.
  \item Predictiva a posteriori: $0{,}5\cdot 0{,}6692+0{,}6\cdot 0{,}3289+0{,}9\cdot 0{,}001927=0{,}3346+0{,}1973+0{,}0017=0{,}5337$.
\end{enumerate}
\resultado \textbf{a posteriori $(0{,}669;\ 0{,}329;\ 0{,}002)$; la probabilidad de cara siguiente pasa de $0{,}667$ a $0{,}534$.}
\comprobacion $0{,}6692+0{,}3289+0{,}0019=1{,}000$. Una sola cara en cinco lanzamientos hace casi imposible la moneda de $p=0{,}9$.
\end{solucion}

\begin{solucion}{eje:4-sesgo}
\datos $X_1,\dots,X_n$ independientes, media $\mu$, varianza $\sigma^{2}$.
\planteamiento escribir $X_i-\Xbar_n=(X_i-\mu)-(\Xbar_n-\mu)$, desarrollar y tomar esperanzas.
\pasos
\begin{enumerate}[label=\arabic*.]
  \item $(X_i-\Xbar_n)^{2}=(X_i-\mu)^{2}-2(X_i-\mu)(\Xbar_n-\mu)+(\Xbar_n-\mu)^{2}$.
  \item Como $\sum_i(X_i-\mu)=n\Xbar_n-n\mu$, el término cruzado sumado vale $-2(\Xbar_n-\mu)\cdot n(\Xbar_n-\mu)=-2n(\Xbar_n-\mu)^{2}$.
  \item El último término, sumado $n$ veces, da $n(\Xbar_n-\mu)^{2}$.
  \item Luego $\sum_i(X_i-\Xbar_n)^{2}=\sum_i(X_i-\mu)^{2}-n(\Xbar_n-\mu)^{2}$.
  \item $\E\bigl[\sum_i(X_i-\mu)^{2}\bigr]=n\sigma^{2}$ por la definición de varianza.
  \item $\E\bigl[n(\Xbar_n-\mu)^{2}\bigr]=n\Var(\Xbar_n)=n\cdot\sigma^{2}/n=\sigma^{2}$ por el \cref{cor:media}.
  \item $\E\bigl[\sum_i(X_i-\Xbar_n)^{2}\bigr]=(n-1)\sigma^{2}$; dividiendo por $n$, $\E[\hat\sigma^{2}]=\frac{n-1}{n}\sigma^{2}$, y dividiendo por $n-1$, $\E[s^{2}]=\sigma^{2}$.
\end{enumerate}
\resultado \textbf{$\E[\hat\sigma^{2}]=(n-1)\sigma^{2}/n$; dividir por $n-1$ elimina el sesgo.} $\blacksquare$
\comprobacion con $n=1$, $\sum_i(X_i-\Xbar_n)^{2}=0$ siempre, y la fórmula da $(1-1)\sigma^{2}=0$. Con $n=5$ y $\sigma^{2}=4$, $\E[\hat\sigma^{2}]=3{,}2$.
\end{solucion}

\begin{solucion}{eje:4-beta}
\datos a priori proporcional a $p^{\alpha-1}(1-p)^{\alpha-1}$; $k$ éxitos en $n$ ensayos; la media de $\BetaD(a,b)$ es $a/(a+b)$.
\planteamiento densidad a posteriori $\propto$ verosimilitud $\times$ a priori, e identificación de la familia.
\pasos
\begin{enumerate}[label=\arabic*.]
  \item Verosimilitud: $\PP(\text{datos}\mid p)=\binom nk p^{k}(1-p)^{n-k}\propto p^{k}(1-p)^{n-k}$.
  \item Producto: $p^{k}(1-p)^{n-k}\cdot p^{\alpha-1}(1-p)^{\alpha-1}=p^{k+\alpha-1}(1-p)^{n-k+\alpha-1}$.
  \item Es el núcleo de una $\BetaD(a,b)$ con $a=k+\alpha$ y $b=n-k+\alpha$; la constante de normalización queda determinada por la condición de integrar~1.
  \item Media: $\dfrac{a}{a+b}=\dfrac{k+\alpha}{(k+\alpha)+(n-k+\alpha)}=\dfrac{k+\alpha}{n+2\alpha}$.
\end{enumerate}
\resultado \textbf{media a posteriori $(k+\alpha)/(n+2\alpha)$; con $\alpha=1$ es el estimador de Laplace, $(k+1)/(n+2)$.} $\blacksquare$
\comprobacion con $n=0$ se obtiene $\alpha/(2\alpha)=1/2$, la media a priori; con $k=7$, $n=10$ y $\alpha=1$, $8/12=2/3$, como en el \cref{eje:4-moneda}.
\end{solucion}

% ---------------------------------------------------------------------
\section{Tema 5}

\begin{solucion}{eje:5-lote}
\datos el tamaño del minilote pasa de $B$ a $4B$.
\planteamiento \cref{prop:minilote}: error típico $\sigma/\sqrt B$; coste proporcional a $B$.
\pasos
\begin{enumerate}[label=\arabic*.]
  \item $\sigma/\sqrt{4B}=\sigma/(2\sqrt B)$: el error típico se divide por~2.
  \item Se procesan $4B$ ejemplos: el coste por paso se multiplica por~4.
\end{enumerate}
\resultado \textbf{el error típico se reduce a la mitad y el coste por paso se cuadruplica.}
\comprobacion con $\sigma=2$: $B=256$ da $0{,}125$ y $B=1024$ da $0{,}0625$.
\end{solucion}

\begin{solucion}{eje:5-vf}
\planteamiento contrastar cada afirmación con la \cref{prop:suma}, la \cref{prop:sesgo}, el \cref{teo:tcl} y la softmax con temperatura del módulo~1.
\pasos
\begin{enumerate}[label=(\alph*)]
  \item \textbf{Falsa.} $\Var(X+Y)=\Var(X)+\Var(Y)+2\Cov(X,Y)$. Con $Y=X$ y $\Var(X)>0$, $\Var(2X)=4\Var(X)\ne 2\Var(X)$.
  \item \textbf{Falsa.} $\E[\hat\sigma^{2}]=(n-1)\sigma^{2}/n<\sigma^{2}$.
  \item \textbf{Falsa.} El teorema exige independencia, igual distribución y varianza finita; la normalidad es la conclusión.
  \item \textbf{Verdadera} si el logit máximo es único: al dividir los logits por $T\to 0$, la diferencia entre el máximo y los demás crece sin límite y la softmax tiende al vector que vale~1 en el máximo.
\end{enumerate}
\resultado \textbf{(a) F; (b) F; (c) F; (d) V.}
\end{solucion}

\begin{solucion}{eje:5-pi}
\datos $n=10^{4}$; $p=\pi/4\approx 0{,}7854$; $p(1-p)=0{,}1685$; valor observado $3{,}19$.
\planteamiento \cref{ej:pi}: $\sigma_{\hat\pi}=4\sqrt{p(1-p)/n}$; para (c), expresar el error en errores típicos y usar el \cref{cor:aprox}.
\pasos
\begin{enumerate}[label=\arabic*.]
  \item (a) $\sigma_{\hat\pi}=4\sqrt{0{,}1685/10^{4}}=4\cdot 0{,}004105=0{,}01642$.
  \item (b) $4\sqrt{0{,}1685/n}=0{,}001\Rightarrow n=16\cdot 0{,}1685/10^{-6}=2{,}696\cdot 10^{6}$.
  \item (c) Error observado: $3{,}19-3{,}1416=0{,}0484$, es decir, $0{,}0484/0{,}01642=2{,}95$ errores típicos.
  \item $\PP(\abs Z\ge 2{,}95)=2(1-\Phi(2{,}95))\approx 0{,}0032$.
\end{enumerate}
\resultado \textbf{(a) $0{,}0164$; (b) unos $2{,}7$ millones de puntos; (c) un error de unos $2{,}9$ errores típicos, con probabilidad $\approx 0{,}3\,\%$: es sospechoso y conviene revisar el código.}
\comprobacion bajar el error de $0{,}0164$ a $0{,}001$ es dividirlo por $16{,}4$; $n$ debe multiplicarse por $16{,}4^{2}\approx 270$, y $10^{4}\cdot 270=2{,}7\cdot 10^{6}$.
\end{solucion}

\begin{solucion}{eje:5-chebyshev}
\datos $p=0{,}5$; $n=100$; $\varepsilon=0{,}1$.
\planteamiento $\hat p$ es la media de $n$ variables $\Ber(p)$: $\Var(\hat p)=p(1-p)/n$ (\cref{cor:media}); \cref{teo:chebyshev}; \cref{cor:aprox}.
\pasos
\begin{enumerate}[label=\arabic*.]
  \item (a) $\Var(\hat p)=0{,}25/100=0{,}0025$; error típico $\sqrt{0{,}0025}=0{,}05$. La varianza escala como $1/n$ y el error típico como $1/\sqrt n$.
  \item (b) $\PP(\abs{\hat p-0{,}5}\ge 0{,}1)\le 0{,}0025/0{,}01=0{,}25$.
  \item (c) $0{,}1$ son $0{,}1/0{,}05=2$ errores típicos; $2(1-\Phi(2))=2(1-0{,}9772)=0{,}0455$.
\end{enumerate}
\resultado \textbf{$\Var(\hat p)=0{,}0025$ y error típico $0{,}05$; Chebyshev da $\le 0{,}25$ y el teorema central del límite, $\approx 0{,}046$.}
\comprobacion Chebyshev es unas $5{,}5$ veces más pesimista: vale para cualquier distribución con esa varianza, mientras que el teorema central del límite usa la forma aproximadamente normal de $\hat p$.
\end{solucion}

\begin{solucion}{eje:5-comparar}
\datos $\hat p_1=0{,}82$, $\hat p_2=0{,}80$, $n=500$ cada uno, independientes.
\planteamiento error típico de cada proporción $\sqrt{\hat p(1-\hat p)/n}$; el de la diferencia de dos independientes, $\sqrt{\sigma_1^{2}+\sigma_2^{2}}$ (\cref{prop:suma}); $z=$ diferencia$/$error típico.
\pasos
\begin{enumerate}[label=\arabic*.]
  \item $\sigma_1^{2}=0{,}82\cdot 0{,}18/500=0{,}1476/500=2{,}952\cdot 10^{-4}$.
  \item $\sigma_2^{2}=0{,}8\cdot 0{,}2/500=0{,}16/500=3{,}2\cdot 10^{-4}$.
  \item $\sigma_{\mathrm{dif}}=\sqrt{6{,}152\cdot 10^{-4}}=0{,}02480$.
  \item (a) $z=0{,}02/0{,}02480=0{,}806<1{,}96$: no significativa.
  \item (b) Con $n$ preguntas por modelo, $\sigma_{\mathrm{dif}}=\sqrt{0{,}3076/n}$. Se impone $1{,}96\sqrt{0{,}3076/n}=0{,}02$.
  \item $n=1{,}96^{2}\cdot 0{,}3076/0{,}02^{2}=3{,}8416\cdot 0{,}3076/0{,}0004=2954{,}2$.
\end{enumerate}
\resultado \textbf{(a) $z\approx 0{,}81$: diferencia no significativa; (b) unas $2955$ preguntas por modelo.}
\comprobacion $n$ crece en el factor $(1{,}96/0{,}806)^{2}=5{,}91$, y $500\cdot 5{,}91\approx 2954$.
\end{solucion}

\begin{solucion}{eje:5-n}
\datos frecuencia $\hat p$ de $n$ ensayos $\Ber(p)$, $p$ desconocido; $\varepsilon=0{,}1$; probabilidad de error admitida $0{,}05$.
\planteamiento (a) completar el cuadrado; (b) cota del \cref{teo:lgn} con la peor varianza; (c) \cref{cor:aprox} con $z=1{,}96$.
\pasos
\begin{enumerate}[label=\arabic*.]
  \item (a) $p(1-p)=\frac14-\bigl(p-\frac12\bigr)^{2}\le\frac14$, con igualdad en $p=1/2$.
  \item (b) $\PP(\abs{\hat p-p}\ge 0{,}1)\le\dfrac{p(1-p)}{n\cdot 0{,}01}\le\dfrac{0{,}25}{0{,}01\,n}=\dfrac{25}{n}$.
  \item Se exige $25/n\le 0{,}05$, es decir, $n\ge 500$.
  \item (c) Con el teorema central del límite y la peor varianza, se pide $1{,}96\cdot 0{,}5/\sqrt n\le 0{,}1$, es decir, $\sqrt n\ge 9{,}8$ y $n\ge 96{,}04$: bastan $97$ ensayos.
\end{enumerate}
\resultado \textbf{$n=500$ garantiza el error pedido para cualquier $p$; el teorema central del límite sugiere que bastan unos 97.} La cota de Chebyshev es rigurosa para todo $n$; la aproximación normal es más ajustada, pero solo asintótica. $\blacksquare$
\comprobacion con $n=500$ y $p=1/2$, el error típico es $\sqrt{0{,}25/500}=0{,}02236$, y $0{,}1$ son $4{,}47$ errores típicos: la probabilidad real de error es muy inferior a $0{,}05$, lo que confirma el carácter pesimista de la cota.
\end{solucion}
